\chapter*{Acrónimos y símbolos}
\addcontentsline{toc}{chapter}{Acrónimos y símbolos}
\thispagestyle{plain}

\section*{Acrónimos}
\begin{longtable}{@{}L{0.18\textwidth}L{0.76\textwidth}@{}}
\toprule
\textbf{Sigla} & \textbf{Significado} \\
\midrule
\endhead
A/D, ADC & Convertidor analógico-digital (\emph{Analog-to-Digital Converter}) \\
API      & Interfaz de programación de aplicaciones \\
CDC      & Clase de dispositivo USB de comunicaciones (puerto serie virtual) \\
CRC      & Código de redundancia cíclica (\emph{Cyclic Redundancy Check}) \\
CSV      & Valores separados por comas (\emph{Comma-Separated Values}) \\
DE / RE  & Habilitación del emisor / receptor de un transceptor RS-485 \\
DI / DO  & Entrada digital / salida digital \\
EEPROM   & Memoria no volátil borrable eléctricamente \\
ESD      & Descarga electrostática (\emph{Electrostatic Discharge}) \\
FC       & Final de carrera \\
GPIO     & Entrada/salida de propósito general \\
GUI      & Interfaz gráfica de usuario \\
HTTP     & Protocolo de transferencia de hipertexto \\
I\textsuperscript{2}C & Bus serie síncrono de dos hilos (\emph{Inter-Integrated Circuit}) \\
LC       & Célula de carga (\emph{Load Cell}); LC1 y LC2 en el firmware \\
MISR     & Registro de firma de la EEPROM del ZSC31014 (\emph{Multiple-Input Signature Register}) \\
NC       & Contacto normalmente cerrado \\
ODS      & Objetivos de Desarrollo Sostenible \\
PCB      & Placa de circuito impreso (\emph{Printed Circuit Board}) \\
PIO      & E/S programable del RP2350 (\emph{Programmable I/O}) \\
POR      & Reinicio por encendido (\emph{Power-On Reset}) \\
RPM, rpm & Revoluciones por minuto \\
RS-485   & Estándar TIA/EIA-485 de transmisión diferencial multipunto \\
RTU      & Modo binario de Modbus sobre línea serie (\emph{Remote Terminal Unit}) \\
SMD      & Componente de montaje superficial \\
STO      & Desconexión segura del par (\emph{Safe Torque Off}) \\
TFG      & Trabajo Fin de Grado \\
UART     & Transmisor-receptor asíncrono universal \\
USB      & Bus serie universal \\
VRD      & Tambor de radio variable (\emph{Variable Radius Drum}) \\
WS       & WebSocket \\
\bottomrule
\end{longtable}

\section*{Símbolos principales}
\begin{longtable}{@{}L{0.17\textwidth}L{0.55\textwidth}L{0.16\textwidth}@{}}
\toprule
\textbf{Símbolo} & \textbf{Descripción} & \textbf{Unidad} \\
\midrule
\endhead
$L_2$     & Altura del punto de salida del cable sobre la articulación O & mm \\
$L_3$     & Distancia de O al punto de anclaje del cable en la barra & mm \\
$L_4$     & Distancia del anclaje del cable al punto de suspensión de la masa & mm \\
$m$       & Masa suspendida & kg \\
$g$       & Aceleración de la gravedad & \si{\metre\per\second\squared} \\
$\theta$  & Ángulo de la barra respecto a la horizontal (positivo hacia arriba) & ° \\
$c$       & Longitud libre de cable entre la salida y el anclaje & mm \\
$T$       & Tensión del cable & N \\
$K$       & Constante de proporcionalidad $T = Kc$ & N/mm \\
$\varphi$ & Giro del tambor (caracola) medido desde el origen del recorrido & rad \\
$r(\varphi)$ & Radio instantáneo de enrollamiento del tambor & mm \\
$s(\varphi)$ & Cable pagado por el tambor, $s=\int r\,\mathrm{d}\varphi$ & mm \\
$\tau$    & Par en el eje del tambor & N·m \\
$N$       & Reducción motor:tambor (vueltas de motor por vuelta de tambor) & -- \\
$\tau_g$, $\tau_f$ & Componentes gravitatoria y de pérdidas del par medido & \ua \\
$p$       & Posición integrada del eje del motor & rev \\
\ua       & Unidades internas del accionamiento (registro 16387) & -- \\
\bottomrule
\end{longtable}
